\section{一般图最大匹配}\label{sec:general-matching}

无向一般图的最大匹配。奇数长度的环无法用二分图方法处理，实现中使用 Edmonds Blossom 缩花。

\subsection{接口}

\apitable{
    \api{GeneralMatching(usize n)}{建 $n$ 个点、零条边的无向图，编号 $[0,n)$}
    \api{usize size()}{点数}
    \api{void add_edge(usize u, usize v)}{加入无向边；内部双向存储，端点应不同}
    \api{std::pair<usize, std::vector<usize>> max_matching()}{返回 \code{{size, match}}（匹配数与匹配表）；\code{match[u]} 为 $u$ 的匹配点，未匹配时为 \code{npos}}
    \api{npos}{表示「未匹配」的哨兵，等于 \code{usize(-1)}}
}

\subsection{约定}

每条无向边只需 \code{add_edge} 一次，内部会同时加入两个方向；重复添加不影响正确性，但会增大常数。端点必须不同，自环在匹配中没有意义，实现也未专门处理。返回 \code{{size, match}}，\code{match} 是对称的，未匹配为 \code{npos}；每次调用都从零开始求解，状态不跨调用保留。

\fullsource{general_matching.hpp}
